%% ============================================================
\section{Preliminaries}\label{sec:preliminaries}
%% ============================================================

\tightparagraph{Loop verification.}
A task contains a loop \(\mathcal L=(\mathsf{Init},B,S)\) and targets
\(\mathcal G=\{(\ell,g_\ell)\}\), where \(\mathsf{Init}\) describes loop entry,
\(B\) is the guard, \(S\) is the body, and \(g_\ell\) occurs at program
location \(\ell\).  In this paper, \emph{valid invariant} means inductive:
\[
  \mathsf{Init}\Rightarrow I,
  \qquad
  \{I\wedge B\}\,S\,\{I\},
\]
which implies every reachable loop-head state satisfies \(I\)
~\citep{floyd1967assigning,hoare1969axiomatic}.  A finite clause set \(C\)
composes as \(I(C)=\bigwedge_{c\in C}c\) and is \emph{jointly inductive} when
\(I(C)\) is valid.  It is target-sufficient when it proves every \(g_\ell\) at
its original location; for an immediate post-loop target with only
guard-false exits and a side-effect-free guard, this reduces to
\(I(C)\wedge\neg B\Rightarrow g\).

\tightparagraph{Targets at different locations.}
Let \(\Sigma\) be the state space and \(S\subseteq\Sigma\times\Sigma\)
the relation for one normally completing loop step.  Write
\(S^B_\ell\) for the path-sensitive body prefix reaching location \(\ell\),
and \(X_\ell\) for paths from a loop head through an exit to that location,
including guard-false and \texttt{break} exits.  Target
sufficiency requires the corresponding obligation for every \((\ell,g)\in\mathcal G\):
\begin{align}
 I(\sigma)\land B(\sigma)\land S^B_\ell(\sigma,\sigma')
   &\Rightarrow g(\sigma') &&\text{(body target)},\nonumber\\
 I(\sigma)\land X_\ell(\sigma,\sigma')
   &\Rightarrow g(\sigma') &&\text{(post-loop target)}.
 \label{eq:main-target-obligations}
\end{align}
Both obligations quantify universally over \(\sigma,\sigma'\).
The relations include feasible-path conditions and executable side effects;
\(X_\ell\) includes body prefixes leading to \texttt{break}, not just
\(\neg B\); Appendix~\ref{app:formal-loop-verification} defines it explicitly.
Thus joint inductiveness and target sufficiency are separate requirements:
an inductive conjunction need not prove the hidden targets, and agreement
with finitely many observed states does not establish induction.

\tightparagraph{Target-hidden verification.}
Generation and intermediate scores use \(\mathcal L\), not \(\mathcal G\).
Targets and target-revealing text are masked, while preconditions and executable
loop semantics remain visible.  The invariant is fixed before the original
targets are restored for final verification.  We study scalar-integer C
programs with one loop; Appendix~\ref{app:scope-example} gives the scope and an
example.
